\subsection{Approximation Methods}

\begin{theorem} \label{thm:weierstrass-polynomial-approximation}\label{thm:power-approx-fourier-main-source-todo} \hlt{Weierstrass Polynomial Approximation}\\
If $f:[a,b]\to\C$ is continuous, then polynomials $P_n$ satisfy $P_n\rightrightarrows f$ on $[a,b]$. If $f$ is real-valued, take $P_n$ real.
\begin{proof}
By an affine change of variables and subtraction of the endpoint interpolant, assume $[a,b]=[0,1]$ and $f(0)=f(1)=0$; extend $f$ by $0$ outside $[0,1]$. For $n\geq1$, let $Q_n(x)\defeq c_n(1-x^2)^n$ and $\int_{-1}^{1}Q_n(x)\,dx=1$.\\
Source estimates give $c_n<\sqrt n$ and $Q_n(t)\leq\sqrt n(1-\delta^2)^n$ for $\delta\leq\abs{t}\leq1$.\\
Define $P_n(x)\defeq\int_{-1}^{1}f(x+t)Q_n(t)\,dt$, $(0\leq x\leq1)$. By change of variables, $P_n(x)=\int_0^1 f(u)Q_n(u-x)\,du$, so $P_n$ is a polynomial in $x$ (real if $f$ is real). Given $\varepsilon>0$, choose $\delta>0$ so $\abs{u-v}<\delta$ implies $\abs{f(u)-f(v)}<\varepsilon/2$, and set $M=\sup\abs{f}$. Since $Q_n\geq0$ and $\int Q_n=1$, $\abs{P_n(x)-f(x)} \leq \int_{-1}^{1}\abs{f(x+t)-f(x)}Q_n(t)\,dt \leq \frac{\varepsilon}{2}+4M\sqrt n(1-\delta^2)^n$.\\
The last term tends to $0$, so $P_n\rightrightarrows f$ on $[0,1]$. Undoing the reductions proves the theorem.\\
\hlt{Structural Commentary}: The kernels $Q_n$ form polynomial approximate identities.\\
\hlt{Key Insight}: Uniform continuity controls the near part, and concentration of $Q_n$ kills the far part.
\end{proof}
\end{theorem}

\begin{figure}[H]
  \centering
  \begin{tikzpicture}[
  >=Stealth,
  every node/.style={font=\scriptsize},
  x=2.55cm,
  y=0.88cm,
  axis/.style={->, black!55, line width=0.45pt},
  broad/.style={black!55, dashed, line width=0.75pt},
  narrow/.style={black!85, line width=0.9pt},
  window/.style={fill=gray!12, draw=none},
  guide/.style={black!35, dashed, line width=0.45pt}
]
  \draw[window] (-0.32,0) rectangle (0.32,2.75);

  \draw[axis] (-1.16,0) -- (1.16,0) node[right] {$t$};
  \draw[axis] (0,-0.15) -- (0,2.86);

  \draw[broad, domain=-1:1, samples=150]
    plot (\x,{1.25*(1-\x*\x)^2});
  \draw[narrow, domain=-1:1, samples=180]
    plot (\x,{2.48*(1-\x*\x)^8});

  \node at (0.66,0.88) {$Q_m$};
  \node at (0.34,1.75) {$Q_n$};

  \draw[guide] (-0.32,0) -- (-0.32,2.55);
  \draw[guide] (0.32,0) -- (0.32,2.55);
  \draw[<->, black!55] (-0.32,-0.22) -- node[below] {$2\delta$} (0.32,-0.22);

  \node[below] at (-1,0) {$-1$};
  \node[below] at (1,0) {$1$};
\end{tikzpicture}

  \caption{The normalized kernels $Q_n(t)=c_n(1-t^2)^n$ concentrate near $0$ while retaining integral $1$.}
  \label{fig:polynomial-bump-kernel}
\end{figure}

\begin{figure}[H]
  \centering
  \begin{tikzpicture}[
  >=Stealth,
  x=4.8cm,
  y=1.75cm,
  every node/.style={font=\scriptsize},
  axis/.style={->, gray!70, line width=0.45pt},
  curve/.style={blue!70!black, line width=0.9pt},
  coarse/.style={orange!82!black, line width=0.75pt},
  fine/.style={teal!70!black, line width=0.65pt},
  point/.style={circle, fill=black, inner sep=1.05pt},
  note/.style={fill=white, inner sep=0.8pt}
]
  \draw[axis] (-0.04,0) -- (1.06,0) node[right] {$x$};
  \draw[axis] (0,-0.06) -- (0,1.08) node[above] {$y$};

  \draw[curve, domain=0:1, variable=\x, samples=140]
    plot (\x,{sqrt(1-\x)});
  \draw[coarse] (0,1) -- (0.5,0.7071) -- (1,0);
  \draw[fine] (0,1) -- (0.125,0.9354) -- (0.25,0.8660) -- (0.375,0.7906)
    -- (0.5,0.7071) -- (0.625,0.6124) -- (0.75,0.5) -- (0.875,0.3536) -- (1,0);

  \foreach \x/\y in {0/1,0.5/0.7071,1/0} {
    \node[point] at (\x,\y) {};
  }

  \node[blue!70!black, note] at (0.30,0.93) {$f$};
  \node[orange!82!black, note] at (0.64,0.56) {$P_2$};
  \node[teal!70!black, note] at (0.82,0.42) {$P_8$};
  \node[below] at (0,0) {$0$};
  \node[below] at (1,0) {$1$};
\end{tikzpicture}

  \caption{Finer piecewise-linear interpolants approach $f(x)=\sqrt{1-x}$ uniformly on $[0,1]$.}
  \label{fig:polygonal-approximation-sqrt-one-minus-x}
\end{figure}

\begin{corollary} \label{cor:absolute-value-polynomial-approximation} \hlt{Polynomial Approximation of $\abs{x}$}\\
For every interval $[-a,a]$, there are real polynomials $P_n$ such that $P_n(0)=0$ and $P_n\rightrightarrows \abs{x}$ on $[-a,a]$.
\begin{proof}
By Theorem \ref{thm:weierstrass-polynomial-approximation}, choose real polynomials $P_n^\ast$ converging uniformly to $\abs{x}$ on $[-a,a]$.\\
Since $P_n^\ast(0)\to0$, the polynomials $P_n(x)\defeq P_n^\ast(x)-P_n^\ast(0)$ have $P_n(0)=0$ and still converge uniformly to $\abs{x}$.\\
\hlt{Structural Commentary}: Normalize the approximants at the base point.\\
\hlt{Key Insight}: Subtracting the value at $0$ fixes the normalization without changing the uniform limit.
\end{proof}
\end{corollary}

\begin{definition} \label{def:algebra-uniform-closure-functions} \hlt{Algebras and Uniform Closure}\\
Let $\mathcal{A}$ be a family of complex-valued functions on a set $E$.
\begin{enumerate}[label=(\roman*)]
\setlength{\itemsep}{0pt}
\item $\mathcal{A}$ is an \hlt{algebra} if it is closed under $+$, products, and scalar multiplication (real scalars in the real case).
\item $\mathcal{A}$ is \hlt{uniformly closed} if $f_n\in\mathcal{A}$ and $f_n\rightrightarrows f$ on $E$ imply $f\in\mathcal{A}$.
\item The \hlt{uniform closure} of $\mathcal{A}$ is $\{f:f_n\in\mathcal{A},\ f_n\rightrightarrows f\text{ on }E\}$.
\end{enumerate}
\end{definition}

\begin{theorem} \label{thm:uniform-closure-algebra} \hlt{Uniform Closures Are Closed Algebras}\\
The uniform closure of an algebra of bounded functions is a uniformly closed algebra.
\begin{proof}
If $f,g\in\mathcal{B}$, choose $f_n,g_n\in\mathcal{A}$ with $f_n\rightrightarrows f$ and $g_n\rightrightarrows g$.\\
Boundedness and uniform convergence give $f_n+g_n\rightrightarrows f+g$; $f_ng_n\rightrightarrows fg$; and $cf_n\rightrightarrows cf$.
Thus $\mathcal{B}$ is an algebra. If $h_n\in\mathcal{B}$ and $h_n\rightrightarrows h$, approximate each $h_n$ uniformly by some $a_n\in\mathcal{A}$ within $1/n$; then $a_n\rightrightarrows h$, so $h\in\mathcal{B}$.\\
\hlt{Structural Commentary}: Uniform limits preserve algebraic operations for bounded functions.\\
\hlt{Key Insight}: The uniform closure of an algebra has already absorbed all further uniform limits.
\end{proof}
\end{theorem}

\begin{definition} \label{def:separates-points-vanishes-nowhere-self-adjoint} \hlt{Separation, Nonvanishing, and Self-Adjointness}\\
Let $\mathcal{A}$ be a family of functions on $E$.
\begin{enumerate}[label=(\roman*)]
\setlength{\itemsep}{0pt}
\item $\mathcal{A}$ \hlt{separates points} on $E$ if for every distinct $x_1,x_2\in E$ there is $f\in\mathcal{A}$ with $f(x_1)\neq f(x_2)$.
\item $\mathcal{A}$ \hlt{vanishes at no point} of $E$ if for every $x\in E$ there is $g\in\mathcal{A}$ with $g(x)\neq0$.
\item A complex algebra $\mathcal{A}$ is \hlt{self-adjoint} if $f\in\mathcal{A}$ implies $\overline{f}\in\mathcal{A}$, where $\overline{f}(x)\defeq\overline{f(x)}$.
\end{enumerate}
\end{definition}

\begin{theorem} \label{thm:two-point-interpolation-in-separating-algebra} \hlt{Two-Point Interpolation in a Separating Algebra}\\
If an algebra $\mathcal{A}$ on $E$ separates points and vanishes nowhere, then for distinct $x_1,x_2\in E$ and scalars $c_1,c_2$, there is $f\in\mathcal{A}$ with $f(x_i)=c_i$, $(i=1,2)$.
\begin{proof}
Choose $g,h,k\in\mathcal{A}$ with $g(x_1)\neq g(x_2)$, $h(x_1)\neq0$, and $k(x_2)\neq0$.\\
Let $u\defeq gk-g(x_1)k$ and $v\defeq gh-g(x_2)h$. Then $u(x_1)=v(x_2)=0$ and $u(x_2)v(x_1)\neq0$.\\
Hence $f\defeq \frac{c_1v}{v(x_1)}+\frac{c_2u}{u(x_2)}$ lies in $\mathcal{A}$ and satisfies $f(x_i)=c_i$.\\
\hlt{Structural Commentary}: Separation and nonvanishing build two functions that isolate the two points.\\
\hlt{Key Insight}: Once two point-isolating functions are in the algebra, interpolation is linear algebra.
\end{proof}
\end{theorem}

\begin{perspective} \label{per:stone-weierstrass-interpolation-linear-algebra} \hlt{Interpolation Is Linear Algebra}\\
The two-point interpolation step can also be written as a $2\times2$ linear system. Its invertibility is exactly the finite-dimensional criterion from Theorem \ref{thm:finite-dimensional-operator-invertible-equivalences}, with the determinant test supplied by Theorem \ref{thm:determinant-basic-laws}. The proof above chooses explicit point-isolating functions instead of reopening the calculation.
\end{perspective}

\begin{theorem} \label{thm:real-stone-weierstrass} \hlt{Stone--Weierstrass Theorem, Real Form}\\
$\mathcal{A}\subset C(K,\R)$ is algebra on compact $K$ that separates points, vanishes nowhere, then uniform closure is $C(K,\R)$.
\begin{proof}
Let $\mathcal{B}$ be the uniform closure of $\mathcal{A}$; by Theorem \ref{thm:uniform-closure-algebra}, $\mathcal{B}$ is a uniformly closed algebra.\\
If $f\in\mathcal{B}$, then $\abs{f}\in\mathcal{B}$: for $a=\norm{f}_{\infty}$, approximate $\abs{y}$ on $[-a,a]$ by polynomials $p_n$ with $p_n(0)=0$ using Corollary \ref{cor:absolute-value-polynomial-approximation}; then $p_n(f)\in\mathcal{B}$ and $p_n(f)\rightrightarrows\abs{f}$.\\
Hence $\max(f,g),\min(f,g)\in\mathcal{B}$ whenever $f,g\in\mathcal{B}$, as $\max(f,g)=\frac{f+g}{2}+\frac{\abs{f-g}}{2}$; $\min(f,g)=\frac{f+g}{2}-\frac{\abs{f-g}}{2}$.\\
Fix $f\in C(K,\R)$, $x\in K$, $\varepsilon>0$. $\forall y\in K$, Theorem \ref{thm:two-point-interpolation-in-separating-algebra} gives $h_y\in\mathcal{B}$, $h_y(x)=f(x)$, $h_y(y)=f(y)$.\\
Compactness gives $y_1,\ldots,y_r$ s.t. $g_x\defeq\max(h_{y_1},\ldots,h_{y_r})$ satisfies $g_x(x)=f(x)$ and $g_x>f-\varepsilon$ on $K$.\\
Compactness gives $x_1,\ldots,x_m$ such that $h\defeq\min(g_{x_1},\ldots,g_{x_m})$ satisfies $\abs{h-f}<\varepsilon$ on $K$.\\
Thus every $f\in C(K,\R)$ is a uniform limit of members of $\mathcal{B}$, so $f\in\mathcal{B}$.\\
\hlt{Structural Commentary}: Local supports plus finite max/min give global approximation.\\
\hlt{Key Insight}: Separation plus compactness gives uniform approximation.
\end{proof}
\end{theorem}

\begin{theorem} \label{thm:complex-stone-weierstrass} \hlt{Stone--Weierstrass Theorem, Complex Form}\\
If $\mathcal{A}\subset C_b(K)$ is self-adjoint, separates points, and vanishes nowhere on compact $K$, then $\overline{\mathcal{A}}=C_b(K)$.
\begin{proof}
Let $\mathcal{A}_{\R}\defeq\{f\in\mathcal{A}:f\text{ real-valued}\}$.\\
If $f=u+iv\in\mathcal{A}$, self-adjointness gives $u=\frac{f+\overline{f}}{2}\in\mathcal{A}_{\R}$ and $v=\frac{f-\overline{f}}{2i}\in\mathcal{A}_{\R}$.\\
Thus $\mathcal{A}_{\R}$ separates points; rescaling a nonzero value and taking real parts shows it vanishes nowhere.\\
By Theorem \ref{thm:real-stone-weierstrass}, $\overline{\mathcal{A}_{\R}}=C(K,\R)$. Approximate the real and imaginary parts separately, then combine.\\
\hlt{Structural Commentary}: Self-adjointness lets the complex algebra expose enough real functions.\\
\hlt{Key Insight}: The complex theorem is the real theorem applied to real and imaginary parts.
\end{proof}
\end{theorem}

\begin{theorem} \label{thm:fundamental-theorem-algebra-exponential-proof} \hlt{Algebraic Completeness of $\C$}\\
Every nonconstant complex polynomial has a complex root.
\begin{proof}
Let $P(z)=\sum_{j=0}^{n}a_jz^j$, $n\geq1$, and normalize $a_n=1$.\\
Since $\abs{P(z)}\to\infty$, Theorem \ref{thm:extreme-value-theorem-compact} gives $z_0$ with $\mu\defeq\abs{P(z_0)}=\min_{\C}\abs{P}$. If $\mu=0$, done.\\
Assume $\mu>0$ and set $Q(w)\defeq P(z_0+w)/P(z_0)$. Then $Q(0)=1$ and $\abs{Q(w)}\geq1$.\\
Write $Q(w)=1+\sum\limits_{j=k}^{n}b_jw^j$, where $b_k\neq0$ and $k$ is minimal. By Theorem \ref{thm:trigonometric-functions-from-exponential}, choose $\theta$ with $b_ke^{ik\theta}=-\abs{b_k}$.\\
Then $Q(re^{i\theta})=1-\abs{b_k}r^k+r^{k+1}H(r)$, with $H$ bounded near $0$, so $\abs{Q(re^{i\theta})}<1$ for small $r>0$, contradiction.\\
\hlt{Structural Commentary}: Minimize, translate, then perturb.\\
\hlt{Key Insight}: Choose the phase that lowers the first nonzero term.
\end{proof}
\end{theorem}

\begin{theorem} \label{thm:gamma-function-basic-theory} \hlt{Gamma Function, Beta Function, and Stirling Skeleton}\\
For $x>0$, $\Gamma(x+1)=x\Gamma(x)$; $\Gamma(n+1)=n!$ for $n\in\N$; $\log\Gamma$ is convex on $(0,\infty)$.\\
Moreover, $\Gamma$ is the unique $f:(0,\infty)\to(0,\infty)$ with $f(1)=1$, $f(x+1)=xf(x)$, and $\log f$ convex.\\
For $x,y>0$, $\mathrm{B}(x,y)\defeq\int_0^1t^{x-1}(1-t)^{y-1}\,dt =\frac{\Gamma(x)\Gamma(y)}{\Gamma(x+y)}$.\\
Consequently, $\Gamma\!\left(\frac12\right)=\sqrt{\pi}$; $\int_{-\infty}^{\infty}e^{-s^2}\,ds=\sqrt{\pi}$; $\Gamma(x)=\frac{2^{x-1}}{\sqrt{\pi}}\Gamma\!\left(\frac{x}{2}\right)\Gamma\!\left(\frac{x+1}{2}\right)$; and $\lim_{x\to\infty}\frac{\Gamma(x+1)}{(x/e)^x\sqrt{2\pi x}}=1$.
\begin{proof}
Integration by parts gives $\Gamma(x+1)=x\Gamma(x)$; $\Gamma(1)=1$ gives $\Gamma(n+1)=n!$.\\
Holder's inequality gives log-convexity. For uniqueness, let $\varphi=\log f$.\\
If $0<x<1$, convexity and $\varphi(x+1)=\varphi(x)+\log x$ give $0\leq \varphi(x)-\log\!\frac{n!\,n^x}{x(x+1)\cdots(x+n)} \leq x\log\!\left(1+\frac1n\right)\to0$.\\
Thus $f$ is uniquely determined on $(0,1)$, hence on $(0,\infty)$.\\
For fixed $y>0$, $\mathrm{B}(1,y)=1/y$, $\mathrm{B}(x+1,y)=\frac{x}{x+y}\mathrm{B}(x,y)$, and Holder gives log-convexity in $x$.\\
Hence $F_y(x)\defeq\frac{\Gamma(x+y)}{\Gamma(y)}\mathrm{B}(x,y)$ satisfies the three uniqueness conditions, so $F_y=\Gamma$.\\
The consequences follow from $t=\sin^2\theta$, $t=s^2$, the duplication identity, and the substitution $t=x(1+u)$ in the asymptotic proof of Stirling's formula.\\
\hlt{Structural Commentary}: Functional equation plus convexity gives uniqueness.\\
\hlt{Key Insight}: Bohr--Mollerup uniqueness turns beta-function identities into gamma identities.
\end{proof}
\end{theorem}

\begin{method} \label{meth:power-approx-fourier-proof-patterns-source-todo}\label{meth:power-series-approximation-proof-patterns} \hlt{Power-Series and Approximation Proof Patterns}
\begin{enumerate}[label=(\roman*)]
\setlength{\itemsep}{0pt}
\item \hlt{Power series}: work on $K_\rho\Subset I_R$; use uniform convergence before termwise limits or derivatives.
\item \hlt{Endpoint theorems}: prove identities inside the radius and pass to endpoints with Abel-type control.
\item \hlt{Approximate identities}: use kernels $Q_n\geq0$, $\int Q_n=1$, with tails $\to0$ off $0$.
\item \hlt{Uniform closure}: show $a_n\in\mathcal{A}$ and $a_n\rightrightarrows f$.
\item \hlt{Lattice closure}: derive $\abs{f}$, then $\max,\min$ from algebraic formulas.
\item \hlt{Stone--Weierstrass}: interpolate locally; use compactness and finite max/min to globalize.
\item \hlt{Gamma function}: use $\Gamma(x+1)=x\Gamma(x)$, $\Gamma(1)=1$, and log-convexity for uniqueness.
\end{enumerate}
\hlt{Key Insight}: Uniform convergence and compactness are what make infinite algebra behave like finite algebra.
\end{method}
